%% ============================================================
%% Appendix B: The Adjoint Formalism — Operators and the Lagrange Identity
%% ============================================================

\chapter{The Adjoint Formalism}
\label{app:adjointformalism}

This appendix develops the abstract operator framework that unifies
the matrix adjoint (Chapter~5) and the ODE adjoint (Chapter~6).
It is intended for students who want a rigorous foundation, not just
the computational recipe.

%%------------------------------------------------------------
\section{Linear Operators and Their Adjoints}
\label{appsec:operators}
%%------------------------------------------------------------

Let $H$ be a Hilbert space with inner product $\langle\cdot,\cdot\rangle$.
A \textbf{linear operator}\index{linear operator} $L : H \to H$ satisfies
$L(\alpha u + \beta v) = \alpha Lu + \beta Lv$.

The \textbf{adjoint operator}\index{adjoint operator} $L^\dagger$ is defined by:
\begin{equation}
  \langle Lu, v \rangle = \langle u, L^\dagger v \rangle
  \qquad\text{for all } u, v \in H.
  \label{eq:adjointdef}
\end{equation}
The adjoint $L^\dagger$ is the unique operator satisfying~\eqref{eq:adjointdef}
(when it exists).

\noindent\textbf{Example~B.1 (Matrix adjoint).}
In $\R^n$ with the Euclidean inner product, $L = A$ (left multiplication).
Then $\langle Au, v\rangle = (Au)^T v = u^T(A^T v) = \langle u, A^T v\rangle$,
so $L^\dagger = A^T$.
This is the matrix transpose: Chapters~5's adjoint equation
$A^T\vec{v} = \vec{c}$ is $L^\dagger\vec{v} = \vec{c}$.

\noindent\textbf{Example~B.2 (Differential operator adjoint).}
In $L^2([0,T])$, let $Lu = du/dt$ with $u(0) = 0$ (zero initial condition).
By integration by parts (Appendix~\ref{appsec:IBP}):
\[
  \langle Lu, v\rangle = \int_0^T \dot{u}v\,dt
  = [uv]_0^T - \int_0^T u\dot{v}\,dt
  = u(T)v(T) - \int_0^T u\dot{v}\,dt.
\]
For the adjoint to satisfy $\langle Lu, v\rangle = \langle u, L^\dagger v\rangle$,
we need $L^\dagger v = -\dot{v}$ with $v(T) = 0$ (zero terminal condition).
So $L^\dagger = -d/dt$ with the boundary condition shifted from $t=0$ to $t=T$.
This is exactly the structure of the adjoint ODE in Chapter~6.

%%------------------------------------------------------------
\section{The Lagrange Identity}
\label{appsec:lagrange}
%%------------------------------------------------------------

For the ODE operator $L = d/dt - F_u$ (the linearization of $\dot{u} = F(u;p)$),
the Lagrange identity\index{Lagrange identity} states:
\begin{equation}
  \frac{d}{dt}\bigl[\lambda^T w\bigr]
  = (L^\dagger\lambda)^T w + \lambda^T (Lw),
  \label{eq:lagrangeID}
\end{equation}
where $w = \partial u/\partial p$ satisfies the FSE $Lw = F_p$ and
$\lambda$ satisfies the adjoint ODE $L^\dagger\lambda = g_u$.

Integrating~\eqref{eq:lagrangeID} over $[0,T]$ and rearranging gives the
sensitivity formula~(6.11) of Chapter~6. This is the rigorous version of
the ``integration by parts'' step that was announced but not fully proved
in the main text.

%%------------------------------------------------------------
\section{Solvability Conditions}
\label{appsec:solvability}
%%------------------------------------------------------------

For the equation $Lu = f$ to have a solution when $L$ is not invertible,
$f$ must be orthogonal to the kernel of $L^\dagger$:

\textbf{SC1 (Fredholm Alternative).}\index{Fredholm alternative}
$Lu = f$ has a solution if and only if $\langle f, v\rangle = 0$
for every $v$ satisfying $L^\dagger v = 0$.

\textbf{SC2 (Uniqueness).}
If $L^\dagger v = 0$ has only the trivial solution $v = 0$,
then $Lu = f$ has a unique solution for every $f$.

In the SA context, SC2 corresponds to the assumption that the Jacobian
at the equilibrium is nonsingular (Theorem~4.1): the FSE has a unique
solution, and the adjoint equation has a unique adjoint variable.
Near a bifurcation, SC2 fails because the Jacobian becomes singular,
which is precisely why SA breaks down (Chapter~8).

%%------------------------------------------------------------
\section{Connection to Chapters 5 and 6}
\label{appsec:connection}
%%------------------------------------------------------------

The following table makes the Chapter~5/6 connection explicit in
the language of this appendix:

\begin{center}
\renewcommand{\arraystretch}{1.4}
\begin{tabular}{lll}
\toprule
Concept & Chapter 5 (matrices) & Chapter 6 (ODEs) \\
\midrule
Space $H$ & $\R^n$ & $L^2([0,T])^n$ \\
Operator $L$ & $A$ (matrix mult.) & $d/dt - J_F$ \\
Adjoint $L^\dagger$ & $A^T$ & $-d/dt - J_F^T$ \\
Adjoint equation & $A^T\vec{v} = \vec{c}$ & $-\dot{\lambda} = J_F^T\lambda - g_u$ \\
Boundary cond. & none (algebraic) & $\lambda(T) = -\nabla h^T$ \\
Sensitivity formula & $\vec{v}^T\vec{r}_j$ & $\int_0^T\lambda^T F_{p_j}\,dt$ \\
Key operation & matrix transpose & integration by parts \\
\bottomrule
\end{tabular}
\end{center}

The formal apparatus of Hilbert spaces and adjoint operators shows that
Chapters~5 and 6 are not two separate methods; they are the same
method applied to two different Hilbert spaces.

%%------------------------------------------------------------
\section{Exercises}
\label{appsec:B_exercises}
%%------------------------------------------------------------

\begin{selfcheck}
Let $A = \begin{pmatrix} 3 & 1 \\ 0 & 2 \end{pmatrix}$
and response $\vec{c} = (1, 0)^T$.
(a) Write the adjoint equation $A^T\vec{v} = \vec{c}$ and solve for $\vec{v}$.
(b) For the right-hand side $\vec{b} = (b_1, b_2)^T$ and solution
$\vec{u} = A^{-1}\vec{b}$, compute $J = \vec{c}^T\vec{u}$ directly.
(c) Verify using the adjoint that $\partial J/\partial b_i = v_i$,
consistent with the sensitivity formula $dJ/d\vec{b} = \vec{v}$.
\end{selfcheck}
\begin{scAnswer}
(a) $A^T = \begin{pmatrix}3&0\\1&2\end{pmatrix}$.
Solving $A^T\vec{v}=(1,0)^T$: $3v_1=1$, $v_1+2v_2=0$, giving $\vec{v}=(1/3,-1/6)^T$.
(b) $A^{-1} = \frac{1}{6}\begin{pmatrix}2&-1\\0&3\end{pmatrix}$.
$\vec{u} = A^{-1}\vec{b} = \frac{1}{6}(2b_1-b_2,\,3b_2)^T$.
$J = u_1 = (2b_1-b_2)/6$.
(c) $\partial J/\partial b_1 = 2/6 = 1/3 = v_1$. \checkmark
$\partial J/\partial b_2 = -1/6 = v_2$. \checkmark
The adjoint gives all sensitivities with one solve.
\end{scAnswer}

\begin{selfcheck}
The Lagrange identity~\eqref{eq:lagrangeID} for the operator
$L = d/dt - F_u$ (the ODE linearization) states:
$\frac{d}{dt}[\lambda^T w] = (L^\dagger\lambda)^T w + \lambda^T(Lw)$.
(a) For the scalar case $L = d/dt - a(t)$ (constant $a$),
write out both sides of the Lagrange identity explicitly.
(b) If $w$ satisfies $Lw = r(t)$ (the FSE with forcing $r$)
and $\lambda$ satisfies $L^\dagger\lambda = 0$, integrate the
Lagrange identity over $[0,T]$ to obtain a formula for
$\lambda(T)w(T) - \lambda(0)w(0)$ in terms of $\int_0^T \lambda r\,dt$.
\end{selfcheck}
\begin{scAnswer}
(a) $L = d/dt - a$, so $L^\dagger = -d/dt - a$ (from integration by parts).
LHS: $d(\lambda w)/dt = \dot\lambda w + \lambda\dot w$.
RHS: $(L^\dagger\lambda)w + \lambda(Lw) = (-\dot\lambda-a\lambda)w + \lambda(\dot w - aw)
= -\dot\lambda w - a\lambda w + \lambda\dot w - a\lambda w$.
LHS $=$ RHS requires $2a\lambda w$ to cancel: checking, $\dot\lambda w + \lambda\dot w
= -\dot\lambda w + \lambda\dot w$, which means $2\dot\lambda w = 0$. This holds when
$L^\dagger\lambda = 0$ means $-\dot\lambda = a\lambda$. So LHS $=$ RHS. \checkmark
(b) Integrating: $[\lambda w]_0^T = \int_0^T(L^\dagger\lambda)^T w\,dt + \int_0^T\lambda(Lw)\,dt
= 0 + \int_0^T\lambda r\,dt$.
So $\lambda(T)w(T) - \lambda(0)w(0) = \int_0^T\lambda(t)r(t)\,dt$.
\end{scAnswer}

\begin{selfcheck}
The Fredholm alternative (Section~\ref{appsec:solvability}) states that
$Lu = f$ has a solution if and only if $f \perp \ker(L^\dagger)$.
For $L = d/dt - a$ on $[0,T]$ with $u(0) = 0$ (zero initial condition),
find $\ker(L^\dagger)$ and state the solvability condition for $Lu = f$
in terms of an integral of $f$.
\end{selfcheck}
\begin{scAnswer}
$L^\dagger = -d/dt - a$ with terminal condition $v(T) = 0$.
$L^\dagger v = 0$ means $-\dot v = av$, so $v(t) = Ce^{-a(t-T)}$.
Applying $v(T) = 0$: $Ce^0 = 0$, so $C = 0$ and $\ker(L^\dagger) = \{0\}$.
Since the kernel is trivial, the Fredholm alternative says $Lu = f$ has a
unique solution for every $f$ (with $u(0) = 0$). This confirms that
when the forward ODE is non-degenerate, the FSE always has a unique solution.
The kernel is nontrivial only when $L$ has a zero eigenvalue, i.e., near a bifurcation.
\end{scAnswer}
